\chapter{ALU Instructions}
\label{ch:alu-instructions}

\section{Overview}

The ALU (Arithmetic Logic Unit) instructions perform arithmetic and logical operations on register operands. All ALU
instructions follow the load/store architecture principle: they operate only on registers, not directly on memory.

ALU instructions are available in three format variants:
\begin{itemize}
	\item \textbf{Immediate}: Operation with 16-bit constant
	\item \textbf{Short Immediate with Shift}: Operation with 8-bit constant and optional shift
	\item \textbf{Register}: Operation between registers with optional shift
\end{itemize}

\section{ALU Instruction Families}

\subsection{Arithmetic Instructions}
\begin{itemize}
	\item \mnemonic{ADD} -- Addition
	\item \mnemonic{ADC} -- Add with Carry
	\item \mnemonic{SUB} -- Subtraction
	\item \mnemonic{SBC} -- Subtract with Carry (Borrow)
\end{itemize}

\subsection{Logical Instructions}
\begin{itemize}
	\item \mnemonic{AND} -- Bitwise AND
	\item \mnemonic{ORR} -- Bitwise OR
	\item \mnemonic{XOR} -- Bitwise XOR (Exclusive OR)
\end{itemize}

\subsection{Data Movement}
\begin{itemize}
	\item \mnemonic{LOAD} -- Load immediate or copy register
\end{itemize}

\subsection{Test and Compare}
\begin{itemize}
	\item \mnemonic{CMP} -- Compare (subtract without saving result)
	\item \mnemonic{TSA} -- Test AND (AND without saving result)
	\item \mnemonic{TSX} -- Test XOR (XOR without saving result)
\end{itemize}

\section{Legal Forms}

ALU instructions are register-only operations. Immediate operands are encoded in the instruction word; memory operands
are not legal for this instruction family.

\begin{table}[H]
	\centering
	\scriptsize
	\begin{tabular}{llll}
		\toprule
		\textbf{Instruction} & \textbf{Syntax}                     & \textbf{Encoding} & \textbf{Notes}                      \\
		\midrule
		ADD                  & \texttt{ADDI rD, \#imm16}           & \opcode{00}       & Writes result                       \\
		ADD                  & \texttt{ADDI rD, \#imm8, shift}     & \opcode{20}       & Writes result                       \\
		ADD                  & \texttt{ADDR rD, rS1, rS2[, shift]} & \opcode{30}       & Writes result                       \\
		ADD                  & \texttt{ADDR rD, rS2[, shift]}      & \opcode{30}       & Writes result (shorthand: rS1 = rD) \\
		ADC                  & \texttt{ADCI rD, \#imm16}           & \opcode{01}       & Writes result                       \\
		ADC                  & \texttt{ADCI rD, \#imm8, shift}     & \opcode{21}       & Writes result                       \\
		ADC                  & \texttt{ADCR rD, rS1, rS2[, shift]} & \opcode{31}       & Writes result                       \\
		SUB                  & \texttt{SUBI rD, \#imm16}           & \opcode{02}       & Writes result                       \\
		SUB                  & \texttt{SUBI rD, \#imm8, shift}     & \opcode{22}       & Writes result                       \\
		SUB                  & \texttt{SUBR rD, rS1, rS2[, shift]} & \opcode{32}       & Writes result                       \\
		SBC                  & \texttt{SBCI rD, \#imm16}           & \opcode{03}       & Writes result                       \\
		SBC                  & \texttt{SBCI rD, \#imm8, shift}     & \opcode{23}       & Writes result                       \\
		SBC                  & \texttt{SBCR rD, rS1, rS2[, shift]} & \opcode{33}       & Writes result                       \\
		AND                  & \texttt{ANDI rD, \#imm16}           & \opcode{04}       & Writes result                       \\
		AND                  & \texttt{ANDI rD, \#imm8, shift}     & \opcode{24}       & Writes result                       \\
		AND                  & \texttt{ANDR rD, rS1, rS2[, shift]} & \opcode{34}       & Writes result                       \\
		ORR                  & \texttt{ORRI rD, \#imm16}           & \opcode{05}       & Writes result                       \\
		ORR                  & \texttt{ORRI rD, \#imm8, shift}     & \opcode{25}       & Writes result                       \\
		ORR                  & \texttt{ORRR rD, rS1, rS2[, shift]} & \opcode{35}       & Writes result                       \\
		XOR                  & \texttt{XORI rD, \#imm16}           & \opcode{06}       & Writes result                       \\
		XOR                  & \texttt{XORI rD, \#imm8, shift}     & \opcode{26}       & Writes result                       \\
		XOR                  & \texttt{XORR rD, rS1, rS2[, shift]} & \opcode{36}       & Writes result                       \\
		LOAD                 & \texttt{LOAD rD, \#imm16}           & \opcode{07}       & Writes result; preserves flags      \\
		LOAD                 & \texttt{LOAD rD, rS}                & \opcode{37}       & Writes result; preserves flags      \\
		CMP                  & \texttt{CMPI rS, \#imm16}           & \opcode{0A}       & Updates flags only                  \\
		CMP                  & \texttt{CMPI rS, \#imm8, shift}     & \opcode{2A}       & Updates flags only                  \\
		CMP                  & \texttt{CMPR rS1, rS2[, shift]}     & \opcode{3A}       & Updates flags only                  \\
		TSA                  & \texttt{TSAI rS, \#imm16}           & \opcode{0C}       & Updates flags only                  \\
		TSA                  & \texttt{TSAI rS, \#imm8, shift}     & \opcode{2C}       & Updates flags only                  \\
		TSA                  & \texttt{TSAR rS1, rS2[, shift]}     & \opcode{3C}       & Updates flags only                  \\
		TSX                  & \texttt{TSXI rS, \#imm16}           & \opcode{0E}       & Updates flags only                  \\
		TSX                  & \texttt{TSXI rS, \#imm8, shift}     & \opcode{2E}       & Updates flags only                  \\
		TSX                  & \texttt{TSXR rS1, rS2[, shift]}     & \opcode{3E}       & Updates flags only                  \\
		\bottomrule
	\end{tabular}
	\caption{Legal ALU Instruction Forms}
	\label{tab:alu-legal-forms}
\end{table}

\noindent\textbf{Register shorthand:} When a register-form ALU instruction omits \texttt{rS1}, the assembler encodes
\texttt{rD} as both the destination and first source operand. For example, \texttt{ADDR r1, r2} is encoded as
\texttt{ADDR r1, r1, r2}. A shift definition may be combined with the shorthand; because the shift always applies to
the left operand, \texttt{ADDR r1, r2, LSL \#2} computes \texttt{(r1 << 2) + r2}, not \texttt{r1 + (r2 << 2)}.

\noindent\textbf{Short-immediate LOAD:} Opcode \texttt{0x27} has no public assembly syntax and is undocumented.

\section{Common ALU Semantics}

\begin{table}[H]
	\centering
	\begin{tabular}{lp{9cm}}
		\toprule
		\textbf{Topic}         & \textbf{Definition}                                                                                                                                                                     \\
		\midrule
		Condition false        & The instruction is skipped and has no side effects. Registers and status flags are preserved.                                                                                           \\
		Immediate format       & The left operand is the named register. The right operand is the encoded 16-bit immediate.                                                                                              \\
		Short immediate format & The left operand is the named register after the encoded shift is applied. The right operand is the encoded 8-bit immediate.                                                            \\
		Register format        & The left operand is \texttt{rS1} after the encoded shift is applied. The right operand is \texttt{rS2}.                                                                                 \\
		Test variants          & \mnemonic{CMP}, \mnemonic{TSA}, and \mnemonic{TSX} compute a result for flag updates but do not write it back.                                                                          \\
		LOAD                   & \mnemonic{LOAD} writes the selected operand to \texttt{rD} and does not update status flags.                                                                                            \\
		Timing                 & ALU instructions complete in 6 cycles, plus any instruction-fetch wait states.                                                                                                          \\
		Exceptions             & Raises no data-memory faults. In protected mode, an ALU instruction whose register field names a privileged register raises a privilege-violation fault at decode; see the notes below. \\
		\bottomrule
	\end{tabular}
	\caption{Common ALU Semantics}
	\label{tab:alu-common-semantics}
\end{table}

\noindent ALU instruction execution does not perform data-memory access and does not raise data bus, data
bus-protection, data alignment, segment-overflow, or invalid-opcode faults. In protected mode, an ALU instruction whose
register field names \reg{sr}, \reg{ah}, \reg{lh}, \reg{ph}, or \reg{sh} raises a privilege-violation fault at decode
and has no other effect. Instruction fetch can still raise the normal fetch faults before an ALU instruction executes.
If trace mode was enabled when the instruction began, an instruction-trace fault is raised after the instruction
commits. Every 6-bit opcode decodes; no processing-unit encoding raises an invalid-opcode fault. Encodings outside the
documented ALU forms execute with architecturally undefined results (see Appendix~\ref{appendix:undocumented}).

\section{Status Flag Updates}

Most ALU instructions update the status register flags:

\begin{description}
	\item[Arithmetic (ADD, ADC, SUB, SBC, CMP)] Update all four flags: N, Z, C, V
	\item[Logical (AND, ORR, XOR, TSA, TSX)] Update N and Z; clear C and V
	\item[Load (LOAD)] Do not update flags
\end{description}

\begin{table}[H]
	\centering
	\begin{tabular}{llllllp{5.5cm}}
		\toprule
		\textbf{Family} & \textbf{N} & \textbf{Z} & \textbf{C} & \textbf{V} & \textbf{Default AF} & \textbf{Meaning}                                                   \\
		\midrule
		ADD             & *          & *          & *          & *          & \texttt{[A]}        & C is carry out; V is signed overflow.                              \\
		ADC             & *          & *          & *          & *          & \texttt{[A]}        & Adds incoming C; C is carry out.                                   \\
		SUB             & *          & *          & *          & *          & \texttt{[A]}        & C is set when a borrow occurs.                                     \\
		SBC             & *          & *          & *          & *          & \texttt{[A]}        & Subtracts incoming C as borrow input; C is set on outgoing borrow. \\
		CMP             & *          & *          & *          & *          & \texttt{[A]}        & Same flags as SUB; result is discarded.                            \\
		AND             & *          & *          & 0          & 0          & \texttt{[A]}        & C and V are cleared.                                               \\
		ORR             & *          & *          & 0          & 0          & \texttt{[A]}        & C and V are cleared.                                               \\
		XOR             & *          & *          & 0          & 0          & \texttt{[A]}        & C and V are cleared.                                               \\
		TSA             & *          & *          & 0          & 0          & \texttt{[A]}        & Same flags as AND; result is discarded.                            \\
		TSX             & *          & *          & 0          & 0          & \texttt{[A]}        & Same flags as XOR; result is discarded.                            \\
		LOAD            & -          & -          & -          & -          & \texttt{[N]}        & Status flags are preserved.                                        \\
		\bottomrule
	\end{tabular}
	\caption{ALU Status Flag Effects}
	\label{tab:alu-status-flag-effects}
\end{table}

See Table~\ref{tab:flag-effect-symbols} for the flag-effect symbol definitions. The table shows the effect when the
instruction uses its default status update source. If \texttt{[N]} is encoded, all status flags are preserved. If
\texttt{[A]} is encoded, flags are updated from the ALU result according to the instruction family. If \texttt{[S]} is
encoded, flags are updated from the shifter result instead of the ALU result; see Chapter~\ref{ch:shift-operations} for
the shifter flag rules. The immediate format has no shift fields, so \texttt{[S]} on an immediate-format instruction
writes a constant zero into the flag byte: N, Z, C, and V are all cleared. Use \texttt{[S]} only with short-immediate
and register forms. Public \mnemonic{LOAD} forms do not support status update override syntax and always preserve the
status flags; an encoding of the \mnemonic{LOAD} opcode with a non-\texttt{None} AF field, which is not reachable from
public assembly syntax, leaves all four status flags architecturally undefined (\texttt{U U U U}).

\subsection{Manual Flag Update Control}

ALU instructions except \mnemonic{LOAD} support manual control over how status flags are updated using the status
override syntax. See Section~\ref{sec:status-override} in Chapter~\ref{ch:shift-operations} for complete details.

\begin{itemize}
	\item \texttt{INSTR[A]} -- Update flags from ALU result (default behavior)
	\item \texttt{INSTR[S]} -- Update flags from shift result instead of ALU
	\item \texttt{INSTR[N]} -- Do not update status flags
	\item AF encoding \texttt{11} (binary) is reserved and has no assembly syntax; the status register is left unchanged.
\end{itemize}

\textbf{Example 13-1:}
\begin{lstlisting}
; Update flags from shift operation
ADDI[S] r1, #0, LSL #3        ; Flags reflect (r1 << 3), not the ALU

; Preserve flags during operation
CMPR r2, r3                   ; Set comparison flags
ADDI[N] r4, #1                ; Increment without changing flags
BRAN|>= @branch_if_r2_gte_r3  ; Use flags from CMPR
\end{lstlisting}

\begin{instructionbox}{SHFT -- Shift with Status Update Meta-Instruction}
	\textbf{Applicability:} SIRC-1, all revisions

	\textbf{Assembles to:} \texttt{ORRI[S] rD, \#0, shift}

	\textbf{Instruction Fields:}
	\begin{itemize}
		\item \textbf{Register (bits 25--22):} \texttt{rD}, the destination and shifted source register; any of the 16
		      registers, \texttt{0x0}--\texttt{0xF}.
		\item \textbf{Immediate (bits 21--14):} fixed at \texttt{0x00} by the meta-instruction.
		\item \textbf{Shift type and count:} \texttt{NUL}, \texttt{LSL}, \texttt{LSR}, \texttt{ASL}, \texttt{ASR},
		      \texttt{RTL}, \texttt{RTR}; an immediate count is \texttt{0}--\texttt{15}, a register count above 15 is
		      architecturally undefined.
		\item \textbf{AF (bits 5--4):} fixed at \texttt{10} (Shift) by the meta-instruction.
		\item \textbf{Condition field (bits 3--0):} any of the 16 condition-code encodings in
		      Table~\ref{tab:condition-code-encoding}.
	\end{itemize}

	\textbf{Syntax:}
	\begin{lstlisting}
SHFT rD, shift                ; rD = shifted rD; flags from shift
\end{lstlisting}

	\textbf{Operands:} \texttt{rD} is both the destination and shifted source register. \texttt{shift} may use any
	standard immediate or register shift count supported by short-immediate ALU instructions.

	\textbf{Operation:}
	\begin{verbatim}
shifted = shift(rD)
rD = shifted OR 0
SR = flags_from_shift_result
\end{verbatim}

	\textbf{Description:}

	Provides concise syntax for shifting a register value while taking status flags from the shifter result. It is not a
	separate CPU operation. The zero OR operation preserves the shifted source value, while \texttt{[S]} makes the status
	flags reflect the shift result instead of the OR result.

	\textbf{Flags:} Updated from the shifter result by default.

	\textbf{Write-back:} Inherits \mnemonic{ORRI} write-back to \texttt{rD}.

	\textbf{Exceptions:} Does not access memory and raises no instruction-specific faults.

	\textbf{Condition field:} If the condition is false, no shift result is written and status flags are preserved.

	\textbf{Timing:} Same as \mnemonic{ORRI}: 6 cycles, plus instruction-fetch wait states.

	\textbf{Privilege:} Available in protected mode and supervisor mode, except when the register field names a
	privileged register (\reg{sr}, \reg{ah}, \reg{lh}, \reg{ph}, \reg{sh}); see Chapter~\ref{ch:exceptions}.

	\textbf{Example 13-2:}
	\begin{lstlisting}
; Shift r1 left by 3 and update flags
SHFT r1, LSL #3               ; ORRI[S] r1, #0, LSL #3

; Shift right and test result
SHFT r3, LSR #2               ; r3 = r3 >> 2
BRAN|== @zero_result          ; Branch if shifted result is zero
\end{lstlisting}

	\textbf{Notes:}
	\begin{itemize}
		\item Immediate shift counts are 0--15; the assembler rejects a larger literal. A register shift count above 15 is
		      architecturally undefined.
		\item All six shift types are supported, plus \mnemonic{NUL}: NUL, LSL, LSR, ASL, ASR, RTL, and RTR. \texttt{SHFT rD, NUL
			      \#0} sets the flags from the unshifted value of \texttt{rD}.
	\end{itemize}
\end{instructionbox}

% ============================================================================
\section{ADD -- Addition}

\begin{instructionbox}{ADDI / ADDR -- Add Immediate / Add Register}
	\textbf{Applicability:} SIRC-1, all revisions

	\textbf{Opcodes:} \opcode{00} (Imm), \opcode{20} (Short Imm), \opcode{30} (Reg)

	\textbf{Instruction Fields:}
	\begin{itemize}
		\item \textbf{Register:} destination \texttt{rD} (immediate and short-immediate forms) or \texttt{rD, rS1, rS2}
		      (register form); any of the 16 registers, \texttt{0x0}--\texttt{0xF}.
		\item \textbf{Immediate:} 16-bit value (immediate form, bits 21--6) or 8-bit value zero-extended to 16 bits
		      (short-immediate form, bits 21--14).
		\item \textbf{Shift type and count (short-immediate and register forms only):} \texttt{NUL}, \texttt{LSL},
		      \texttt{LSR}, \texttt{ASL}, \texttt{ASR}, \texttt{RTL}, \texttt{RTR}; an immediate count is
		      \texttt{0}--\texttt{15}; a register count is used in full (16 or more clamps to a full shift for \texttt{LSL}/\texttt{LSR}/\texttt{ASL}/\texttt{ASR}, and is taken modulo 16 for \texttt{RTL}/\texttt{RTR}) -- see Chapter~\ref{ch:shift-operations}.
		\item \textbf{AF (bits 5--4):} \texttt{00} = None, \texttt{01} = Alu (default), \texttt{10} = Shift, \texttt{11} =
		      Reserved (status register unchanged).
		\item \textbf{Condition field (bits 3--0):} any of the 16 condition-code encodings in
		      Table~\ref{tab:condition-code-encoding}.
	\end{itemize}

	\textbf{Syntax:}
	\begin{lstlisting}
ADDI rD, #imm16                ; rD = rD + imm16
ADDI rD, #imm8, shift          ; rD = shift(rD) + imm8
ADDR rD, rS1, rS2              ; rD = rS1 + rS2
ADDR rD, rS1, rS2, shift       ; rD = shift(rS1) + rS2
\end{lstlisting}

	\textbf{Operands:} \texttt{rD} is the destination and first source in immediate forms. Register forms use \texttt{rS1}
	as the first source and \texttt{rS2} as the second source; assembler shorthand may omit \texttt{rS1}. Short-immediate
	and register forms may apply a shift to the first source operand.

	\textbf{Operation:}
	\begin{verbatim}
result = operand1 + operand2
destination = result
SR.N = result[15]
SR.Z = (result == 0)
SR.C = carry_out
SR.V = signed_overflow
\end{verbatim}

	\textbf{Description:}

	Adds two values and stores the result in the destination register. In short immediate format with shift, the
	destination register is shifted before adding the immediate. The carry flag is set if an unsigned overflow occurs
	(carry out of bit 15). The overflow flag is set if a signed overflow occurs (the sign of the result is incorrect for
	the operation).

	\textbf{Flags:} N Z C V

	\textbf{Write-back:} If the condition is true, writes the 16-bit result to \texttt{rD}.

	\textbf{Exceptions:} Does not access memory and raises no instruction-specific faults.

	\textbf{Condition field:} If the condition is false, no result is written and status flags are preserved.

	\textbf{Timing:} 6 cycles, plus instruction-fetch wait states.

	\textbf{Privilege:} Available in protected mode and supervisor mode, except when the register field names a
	privileged register (\reg{sr}, \reg{ah}, \reg{lh}, \reg{ph}, \reg{sh}); see Chapter~\ref{ch:exceptions}.

	\textbf{Example 13-3:}
	\begin{lstlisting}
; Simple addition
ADDI r1, #100              ; r1 = r1 + 100

; Add two registers
ADDR r3, r1, r2            ; r3 = r1 + r2

; Scaled addition (multiply by 3)
ADDR r4, r5, r5, LSL #1    ; r4 = r5 + (r5 * 2) = r5 * 3

; Build large constant
ADDI r1, #1, LSL #12       ; r1 = (r1 << 12) + 1
\end{lstlisting}

	\textbf{Notes:}
	\begin{itemize}
		\item For multi-precision arithmetic, use \mnemonic{ADC} to propagate carry
		\item Adding zero (\texttt{ADDI r1, \#0}) can be used to move values or test flags
		\item Shift operations allow efficient constant multiplication
	\end{itemize}
\end{instructionbox}

% ============================================================================
\section{ADC -- Add with Carry}

\begin{instructionbox}{ADCI / ADCR -- Add with Carry}
	\textbf{Applicability:} SIRC-1, all revisions

	\textbf{Opcodes:} \opcode{01} (Imm), \opcode{21} (Short Imm), \opcode{31} (Reg)

	\textbf{Instruction Fields:}
	\begin{itemize}
		\item \textbf{Register:} destination \texttt{rD} (immediate and short-immediate forms) or \texttt{rD, rS1, rS2}
		      (register form); any of the 16 registers, \texttt{0x0}--\texttt{0xF}.
		\item \textbf{Immediate:} 16-bit value (immediate form, bits 21--6) or 8-bit value zero-extended to 16 bits
		      (short-immediate form, bits 21--14).
		\item \textbf{Shift type and count (short-immediate and register forms only):} \texttt{NUL}, \texttt{LSL},
		      \texttt{LSR}, \texttt{ASL}, \texttt{ASR}, \texttt{RTL}, \texttt{RTR}; an immediate count is
		      \texttt{0}--\texttt{15}; a register count is used in full (16 or more clamps to a full shift for \texttt{LSL}/\texttt{LSR}/\texttt{ASL}/\texttt{ASR}, and is taken modulo 16 for \texttt{RTL}/\texttt{RTR}) -- see Chapter~\ref{ch:shift-operations}.
		\item \textbf{AF (bits 5--4):} \texttt{00} = None, \texttt{01} = Alu (default), \texttt{10} = Shift, \texttt{11} =
		      Reserved (status register unchanged).
		\item \textbf{Condition field (bits 3--0):} any of the 16 condition-code encodings in
		      Table~\ref{tab:condition-code-encoding}.
	\end{itemize}

	\textbf{Syntax:}
	\begin{lstlisting}
ADCI rD, #imm16                ; rD = rD + imm16 + C
ADCI rD, #imm8, shift          ; rD = shift(rD) + imm8 + C
ADCR rD, rS1, rS2              ; rD = rS1 + rS2 + C
ADCR rD, rS1, rS2, shift       ; rD = shift(rS1) + rS2 + C
\end{lstlisting}

	\textbf{Operands:} Same operand forms as \mnemonic{ADD}. The incoming \texttt{C} flag is read as an additional
	least-significant carry input.

	\textbf{Operation:}
	\begin{verbatim}
result = operand1 + operand2 + SR.C
destination = result
SR.N = result[15]
SR.Z = (result == 0)
SR.C = carry_out
SR.V = signed_overflow
\end{verbatim}

	\textbf{Description:}

	Adds two values plus the carry flag and stores the result. In short immediate format with shift, the destination
	register is shifted before adding the immediate. This instruction is used for multi-precision (32-bit, 48-bit, etc.)
	arithmetic to propagate the carry from lower-order additions to higher-order additions.

	\textbf{Flags:} N Z C V

	\textbf{Write-back:} If the condition is true, writes the 16-bit result to \texttt{rD}.

	\textbf{Exceptions:} Does not access memory and raises no instruction-specific faults.

	\textbf{Condition field:} If the condition is false, no result is written and status flags are preserved.

	\textbf{Timing:} 6 cycles, plus instruction-fetch wait states.

	\textbf{Privilege:} Available in protected mode and supervisor mode, except when the register field names a
	privileged register (\reg{sr}, \reg{ah}, \reg{lh}, \reg{ph}, \reg{sh}); see Chapter~\ref{ch:exceptions}.

	\textbf{Example 13-4:}
	\begin{lstlisting}
; 32-bit addition (r1:r2 = r3:r4 + r5:r6)
; Low words
ADDR r2, r4, r6            ; r2 = r4 + r6, set carry

; High words (with carry propagation)
ADCR r1, r3, r5            ; r1 = r3 + r5 + carry
\end{lstlisting}

	\textbf{Notes:}
	\begin{itemize}
		\item Always use after \mnemonic{ADD} for multi-word arithmetic
		\item The carry flag must be in the correct state before \mnemonic{ADC}
		\item Can be used to add 1 conditionally based on carry flag
	\end{itemize}
\end{instructionbox}

% ============================================================================
\section{SUB -- Subtraction}

\begin{instructionbox}{SUBI / SUBR -- Subtract Immediate / Subtract Register}
	\textbf{Applicability:} SIRC-1, all revisions

	\textbf{Opcodes:} \opcode{02} (Imm), \opcode{22} (Short Imm), \opcode{32} (Reg)

	\textbf{Instruction Fields:}
	\begin{itemize}
		\item \textbf{Register:} destination \texttt{rD} (immediate and short-immediate forms) or \texttt{rD, rS1, rS2}
		      (register form); any of the 16 registers, \texttt{0x0}--\texttt{0xF}.
		\item \textbf{Immediate:} 16-bit value (immediate form, bits 21--6) or 8-bit value zero-extended to 16 bits
		      (short-immediate form, bits 21--14).
		\item \textbf{Shift type and count (short-immediate and register forms only):} \texttt{NUL}, \texttt{LSL},
		      \texttt{LSR}, \texttt{ASL}, \texttt{ASR}, \texttt{RTL}, \texttt{RTR}; an immediate count is
		      \texttt{0}--\texttt{15}; a register count is used in full (16 or more clamps to a full shift for \texttt{LSL}/\texttt{LSR}/\texttt{ASL}/\texttt{ASR}, and is taken modulo 16 for \texttt{RTL}/\texttt{RTR}) -- see Chapter~\ref{ch:shift-operations}.
		\item \textbf{AF (bits 5--4):} \texttt{00} = None, \texttt{01} = Alu (default), \texttt{10} = Shift, \texttt{11} =
		      Reserved (status register unchanged).
		\item \textbf{Condition field (bits 3--0):} any of the 16 condition-code encodings in
		      Table~\ref{tab:condition-code-encoding}.
	\end{itemize}

	\textbf{Syntax:}
	\begin{lstlisting}
SUBI rD, #imm16                ; rD = rD - imm16
SUBI rD, #imm8, shift          ; rD = shift(rD) - imm8
SUBR rD, rS1, rS2              ; rD = rS1 - rS2
SUBR rD, rS1, rS2, shift       ; rD = shift(rS1) - rS2
\end{lstlisting}

	\textbf{Operands:} Same operand forms as \mnemonic{ADD}. The second operand is subtracted from the first operand.

	\textbf{Operation:}
	\begin{verbatim}
result = operand1 - operand2
destination = result
SR.N = result[15]
SR.Z = (result == 0)
SR.C = borrow
SR.V = signed_overflow
\end{verbatim}

	\textbf{Description:}

	Subtracts the second operand from the first and stores the result. In short immediate format with shift, the
	destination register is shifted before subtracting the immediate. The carry flag is set if a borrow occurs (operand2 >
	operand1 for unsigned). The overflow flag indicates signed overflow.

	\textbf{Flags:} N Z C V

	\textbf{Write-back:} If the condition is true, writes the 16-bit result to \texttt{rD}.

	\textbf{Exceptions:} Does not access memory and raises no instruction-specific faults.

	\textbf{Condition field:} If the condition is false, no result is written and status flags are preserved.

	\textbf{Timing:} 6 cycles, plus instruction-fetch wait states.

	\textbf{Privilege:} Available in protected mode and supervisor mode, except when the register field names a
	privileged register (\reg{sr}, \reg{ah}, \reg{lh}, \reg{ph}, \reg{sh}); see Chapter~\ref{ch:exceptions}.

	\textbf{Example 13-5:}
	\begin{lstlisting}
; Decrement by 1
SUBI r1, #1                ; r1 = r1 - 1

; Subtract registers
SUBR r3, r1, r2            ; r3 = r1 - r2

; Loop counter
SUBI r7, #1                ; Decrement counter
BRAN|!= @loop              ; Branch if not zero
\end{lstlisting}

	\textbf{Notes:}
	\begin{itemize}
		\item Carry is set when the subtraction requires a borrow
		\item For multi-precision subtraction, use \mnemonic{SBC}
		\item Subtracting a value from itself (\texttt{SUBR r1, r1, r1}) clears the register
	\end{itemize}
\end{instructionbox}

% ============================================================================
\section{SBC -- Subtract with Carry (Borrow)}

\begin{instructionbox}{SBCI / SBCR -- Subtract with Carry}
	\textbf{Applicability:} SIRC-1, all revisions

	\textbf{Opcodes:} \opcode{03} (Imm), \opcode{23} (Short Imm), \opcode{33} (Reg)

	\textbf{Instruction Fields:}
	\begin{itemize}
		\item \textbf{Register:} destination \texttt{rD} (immediate and short-immediate forms) or \texttt{rD, rS1, rS2}
		      (register form); any of the 16 registers, \texttt{0x0}--\texttt{0xF}.
		\item \textbf{Immediate:} 16-bit value (immediate form, bits 21--6) or 8-bit value zero-extended to 16 bits
		      (short-immediate form, bits 21--14).
		\item \textbf{Shift type and count (short-immediate and register forms only):} \texttt{NUL}, \texttt{LSL},
		      \texttt{LSR}, \texttt{ASL}, \texttt{ASR}, \texttt{RTL}, \texttt{RTR}; an immediate count is
		      \texttt{0}--\texttt{15}; a register count is used in full (16 or more clamps to a full shift for \texttt{LSL}/\texttt{LSR}/\texttt{ASL}/\texttt{ASR}, and is taken modulo 16 for \texttt{RTL}/\texttt{RTR}) -- see Chapter~\ref{ch:shift-operations}.
		\item \textbf{AF (bits 5--4):} \texttt{00} = None, \texttt{01} = Alu (default), \texttt{10} = Shift, \texttt{11} =
		      Reserved (status register unchanged).
		\item \textbf{Condition field (bits 3--0):} any of the 16 condition-code encodings in
		      Table~\ref{tab:condition-code-encoding}.
	\end{itemize}

	\textbf{Syntax:}
	\begin{lstlisting}
SBCI rD, #imm16                ; rD = rD - imm16 - C
SBCI rD, #imm8, shift          ; rD = shift(rD) - imm8 - C
SBCR rD, rS1, rS2              ; rD = rS1 - rS2 - C
SBCR rD, rS1, rS2, shift       ; rD = shift(rS1) - rS2 - C
\end{lstlisting}

	\textbf{Operands:} Same operand forms as \mnemonic{SUB}. The incoming \texttt{C} flag is read as the borrow input.

	\textbf{Operation:}
	\begin{verbatim}
result = operand1 - operand2 - SR.C
destination = result
SR.N = result[15]
SR.Z = (result == 0)
SR.C = borrow
SR.V = signed_overflow
\end{verbatim}

	\textbf{Description:}

	Subtracts the second operand and borrow from the first operand. In short immediate format with shift, the destination
	register is shifted before subtracting the immediate. Used for multi-precision subtraction to propagate borrows from
	lower-order to higher-order subtractions.

	\textbf{Flags:} N Z C V

	\textbf{Write-back:} If the condition is true, writes the 16-bit result to \texttt{rD}.

	\textbf{Exceptions:} Does not access memory and raises no instruction-specific faults.

	\textbf{Condition field:} If the condition is false, no result is written and status flags are preserved.

	\textbf{Timing:} 6 cycles, plus instruction-fetch wait states.

	\textbf{Privilege:} Available in protected mode and supervisor mode, except when the register field names a
	privileged register (\reg{sr}, \reg{ah}, \reg{lh}, \reg{ph}, \reg{sh}); see Chapter~\ref{ch:exceptions}.

	\textbf{Example 13-6:}
	\begin{lstlisting}
; 32-bit subtraction (r1:r2 = r3:r4 - r5:r6)
; Low words
SUBR r2, r4, r6            ; r2 = r4 - r6, set/clear carry

; High words (with borrow propagation)
SBCR r1, r3, r5            ; r1 = r3 - r5 - borrow
\end{lstlisting}

	\textbf{Notes:}
	\begin{itemize}
		\item Always use after \mnemonic{SUB} for multi-word subtraction
		\item Carry is treated as the borrow input for multi-word subtraction
	\end{itemize}
\end{instructionbox}

% ============================================================================
\section{AND -- Bitwise AND}

\begin{instructionbox}{ANDI / ANDR -- AND Immediate / AND Register}
	\textbf{Applicability:} SIRC-1, all revisions

	\textbf{Opcodes:} \opcode{04} (Imm), \opcode{24} (Short Imm), \opcode{34} (Reg)

	\textbf{Instruction Fields:}
	\begin{itemize}
		\item \textbf{Register:} destination \texttt{rD} (immediate and short-immediate forms) or \texttt{rD, rS1, rS2}
		      (register form); any of the 16 registers, \texttt{0x0}--\texttt{0xF}.
		\item \textbf{Immediate:} 16-bit value (immediate form, bits 21--6) or 8-bit value zero-extended to 16 bits
		      (short-immediate form, bits 21--14).
		\item \textbf{Shift type and count (short-immediate and register forms only):} \texttt{NUL}, \texttt{LSL},
		      \texttt{LSR}, \texttt{ASL}, \texttt{ASR}, \texttt{RTL}, \texttt{RTR}; an immediate count is
		      \texttt{0}--\texttt{15}; a register count is used in full (16 or more clamps to a full shift for \texttt{LSL}/\texttt{LSR}/\texttt{ASL}/\texttt{ASR}, and is taken modulo 16 for \texttt{RTL}/\texttt{RTR}) -- see Chapter~\ref{ch:shift-operations}.
		\item \textbf{AF (bits 5--4):} \texttt{00} = None, \texttt{01} = Alu (default), \texttt{10} = Shift, \texttt{11} =
		      Reserved (status register unchanged).
		\item \textbf{Condition field (bits 3--0):} any of the 16 condition-code encodings in
		      Table~\ref{tab:condition-code-encoding}.
	\end{itemize}

	\textbf{Syntax:}
	\begin{lstlisting}
ANDI rD, #imm16                ; rD = rD & imm16
ANDI rD, #imm8, shift          ; rD = shift(rD) & imm8
ANDR rD, rS1, rS2              ; rD = rS1 & rS2
ANDR rD, rS1, rS2, shift       ; rD = shift(rS1) & rS2
\end{lstlisting}

	\textbf{Operands:} Same operand forms as \mnemonic{ADD}. Logical operations interpret operands as 16-bit bit
	patterns; short immediates are zero-extended.

	\textbf{Operation:}
	\begin{verbatim}
result = operand1 AND operand2
destination = result
SR.N = result[15]
SR.Z = (result == 0)
SR.C = 0
SR.V = 0
\end{verbatim}

	\textbf{Description:}

	Performs bitwise AND operation. Each bit in the result is 1 only if the corresponding bits in both operands are 1. In
	short immediate format with shift, the destination register is shifted before ANDing with the immediate. Commonly used
	for bit masking and clearing specific bits.

	\textbf{Flags:} N Z (C and V cleared)

	\textbf{Write-back:} If the condition is true, writes the 16-bit result to \texttt{rD}.

	\textbf{Exceptions:} Does not access memory and raises no instruction-specific faults.

	\textbf{Condition field:} If the condition is false, no result is written and status flags are preserved.

	\textbf{Timing:} 6 cycles, plus instruction-fetch wait states.

	\textbf{Privilege:} Available in protected mode and supervisor mode, except when the register field names a
	privileged register (\reg{sr}, \reg{ah}, \reg{lh}, \reg{ph}, \reg{sh}); see Chapter~\ref{ch:exceptions}.

	\textbf{Example 13-7:}
	\begin{lstlisting}
; Mask lower byte
ANDI r1, #0x00FF           ; r1 = r1 & 0xFF

; Clear bit 5
ANDI r2, #0xFFDF           ; r2 = r2 & 0xFFDF (clear bit 5)

; Check if two registers share any set bits
ANDR r3, r1, r2            ; r3 = r1 & r2
CMPI r3, #0
BRAN|!= @common_bits       ; Branch if any bits in common
\end{lstlisting}

	\textbf{Notes:}
	\begin{itemize}
		\item ANDing with 0xFFFF leaves value unchanged
		\item ANDing with 0x0000 clears the register
		\item Use \mnemonic{TSA} to test bits without modifying the register
	\end{itemize}
\end{instructionbox}

% ============================================================================
\section{ORR -- Bitwise OR}

\begin{instructionbox}{ORRI / ORRR -- OR Immediate / OR Register}
	\textbf{Applicability:} SIRC-1, all revisions

	\textbf{Opcodes:} \opcode{05} (Imm), \opcode{25} (Short Imm), \opcode{35} (Reg)

	\textbf{Instruction Fields:}
	\begin{itemize}
		\item \textbf{Register:} destination \texttt{rD} (immediate and short-immediate forms) or \texttt{rD, rS1, rS2}
		      (register form); any of the 16 registers, \texttt{0x0}--\texttt{0xF}.
		\item \textbf{Immediate:} 16-bit value (immediate form, bits 21--6) or 8-bit value zero-extended to 16 bits
		      (short-immediate form, bits 21--14).
		\item \textbf{Shift type and count (short-immediate and register forms only):} \texttt{NUL}, \texttt{LSL},
		      \texttt{LSR}, \texttt{ASL}, \texttt{ASR}, \texttt{RTL}, \texttt{RTR}; an immediate count is
		      \texttt{0}--\texttt{15}; a register count is used in full (16 or more clamps to a full shift for \texttt{LSL}/\texttt{LSR}/\texttt{ASL}/\texttt{ASR}, and is taken modulo 16 for \texttt{RTL}/\texttt{RTR}) -- see Chapter~\ref{ch:shift-operations}.
		\item \textbf{AF (bits 5--4):} \texttt{00} = None, \texttt{01} = Alu (default), \texttt{10} = Shift, \texttt{11} =
		      Reserved (status register unchanged).
		\item \textbf{Condition field (bits 3--0):} any of the 16 condition-code encodings in
		      Table~\ref{tab:condition-code-encoding}.
	\end{itemize}

	\textbf{Syntax:}
	\begin{lstlisting}
ORRI rD, #imm16                ; rD = rD | imm16
ORRI rD, #imm8, shift          ; rD = shift(rD) | imm8
ORRR rD, rS1, rS2              ; rD = rS1 | rS2
ORRR rD, rS1, rS2, shift       ; rD = shift(rS1) | rS2
\end{lstlisting}

	\textbf{Operands:} Same operand forms as \mnemonic{AND}. Logical operations interpret operands as 16-bit bit
	patterns; short immediates are zero-extended.

	\textbf{Operation:}
	\begin{verbatim}
result = operand1 OR operand2
destination = result
SR.N = result[15]
SR.Z = (result == 0)
SR.C = 0
SR.V = 0
\end{verbatim}

	\textbf{Description:}

	Performs bitwise OR operation. Each bit in the result is 1 if either corresponding bit in the operands is 1. In short
	immediate format with shift, the destination register is shifted before ORing with the immediate. Commonly used for
	setting specific bits.

	\textbf{Flags:} N Z (C and V cleared)

	\textbf{Write-back:} If the condition is true, writes the 16-bit result to \texttt{rD}.

	\textbf{Exceptions:} Does not access memory and raises no instruction-specific faults.

	\textbf{Condition field:} If the condition is false, no result is written and status flags are preserved.

	\textbf{Timing:} 6 cycles, plus instruction-fetch wait states.

	\textbf{Privilege:} Available in protected mode and supervisor mode, except when the register field names a
	privileged register (\reg{sr}, \reg{ah}, \reg{lh}, \reg{ph}, \reg{sh}); see Chapter~\ref{ch:exceptions}.

	\textbf{Example 13-8:}
	\begin{lstlisting}
; Set bit 7
ORRI r1, #0x0080           ; r1 = r1 | 0x80

; Set multiple bits
ORRI r2, #0xF000           ; Set high nibble

; Combine two registers
ORRR r3, r1, r2            ; r3 = r1 | r2

; Set bit N (variable)
LOAD r4, #1                ; r4 = 1
SHFT r4, LSL r5            ; r4 = 1 << r5
ORRR r1, r4                ; r1 |= (1 << r5)
\end{lstlisting}

	\textbf{Notes:}
	\begin{itemize}
		\item ORing with 0x0000 leaves value unchanged
		\item ORing with 0xFFFF sets all bits
		\item Use for enabling flag bits or combining bit fields
	\end{itemize}
\end{instructionbox}

% ============================================================================
\section{XOR -- Bitwise Exclusive OR}

\begin{instructionbox}{XORI / XORR -- XOR Immediate / XOR Register}
	\textbf{Applicability:} SIRC-1, all revisions

	\textbf{Opcodes:} \opcode{06} (Imm), \opcode{26} (Short Imm), \opcode{36} (Reg)

	\textbf{Instruction Fields:}
	\begin{itemize}
		\item \textbf{Register:} destination \texttt{rD} (immediate and short-immediate forms) or \texttt{rD, rS1, rS2}
		      (register form); any of the 16 registers, \texttt{0x0}--\texttt{0xF}.
		\item \textbf{Immediate:} 16-bit value (immediate form, bits 21--6) or 8-bit value zero-extended to 16 bits
		      (short-immediate form, bits 21--14).
		\item \textbf{Shift type and count (short-immediate and register forms only):} \texttt{NUL}, \texttt{LSL},
		      \texttt{LSR}, \texttt{ASL}, \texttt{ASR}, \texttt{RTL}, \texttt{RTR}; an immediate count is
		      \texttt{0}--\texttt{15}; a register count is used in full (16 or more clamps to a full shift for \texttt{LSL}/\texttt{LSR}/\texttt{ASL}/\texttt{ASR}, and is taken modulo 16 for \texttt{RTL}/\texttt{RTR}) -- see Chapter~\ref{ch:shift-operations}.
		\item \textbf{AF (bits 5--4):} \texttt{00} = None, \texttt{01} = Alu (default), \texttt{10} = Shift, \texttt{11} =
		      Reserved (status register unchanged).
		\item \textbf{Condition field (bits 3--0):} any of the 16 condition-code encodings in
		      Table~\ref{tab:condition-code-encoding}.
	\end{itemize}

	\textbf{Syntax:}
	\begin{lstlisting}
XORI rD, #imm16                ; rD = rD ^ imm16
XORI rD, #imm8, shift          ; rD = shift(rD) ^ imm8
XORR rD, rS1, rS2              ; rD = rS1 ^ rS2
XORR rD, rS1, rS2, shift       ; rD = shift(rS1) ^ rS2
\end{lstlisting}

	\textbf{Operands:} Same operand forms as \mnemonic{AND}. Logical operations interpret operands as 16-bit bit
	patterns; short immediates are zero-extended.

	\textbf{Operation:}
	\begin{verbatim}
result = operand1 XOR operand2
destination = result
SR.N = result[15]
SR.Z = (result == 0)
SR.C = 0
SR.V = 0
\end{verbatim}

	\textbf{Description:}

	Performs bitwise XOR (exclusive OR) operation. Each bit in the result is 1 if the corresponding bits in the operands
	are different. In short immediate format with shift, the destination register is shifted before XORing with the
	immediate. Used for toggling bits, bitwise comparison, and simple encryption.

	\textbf{Flags:} N Z (C and V cleared)

	\textbf{Write-back:} If the condition is true, writes the 16-bit result to \texttt{rD}.

	\textbf{Exceptions:} Does not access memory and raises no instruction-specific faults.

	\textbf{Condition field:} If the condition is false, no result is written and status flags are preserved.

	\textbf{Timing:} 6 cycles, plus instruction-fetch wait states.

	\textbf{Privilege:} Available in protected mode and supervisor mode, except when the register field names a
	privileged register (\reg{sr}, \reg{ah}, \reg{lh}, \reg{ph}, \reg{sh}); see Chapter~\ref{ch:exceptions}.

	\textbf{Example 13-9:}
	\begin{lstlisting}
; Toggle bit 3
XORI r1, #0x0008           ; r1 = r1 ^ 0x08

; Invert all bits
XORI r2, #0xFFFF           ; r2 = ~r2

; Check if registers are equal
XORR r3, r1, r2            ; r3 = r1 ^ r2
; If r3 == 0, then r1 == r2

; Swap using XOR (no temporary)
XORR r1, r1, r2            ; r1 = r1 ^ r2
XORR r2, r1, r2            ; r2 = r1 ^ r2 = (r1 ^ r2) ^ r2 = r1
XORR r1, r1, r2            ; r1 = (r1 ^ r2) ^ r1 = r2
\end{lstlisting}

	\textbf{Notes:}
	\begin{itemize}
		\item XORing with 0x0000 leaves value unchanged
		\item XORing with 0xFFFF inverts all bits (bitwise NOT)
		\item XORing a value with itself always gives 0
		\item Use \mnemonic{TSX} to check for differences without modifying register
	\end{itemize}
\end{instructionbox}

% ============================================================================
\section{LOAD -- Load Immediate / Move Register}

\begin{instructionbox}{LOAD -- Load/Move}
	\textbf{Applicability:} SIRC-1, all revisions

	\textbf{Opcodes:} \opcode{07} (Imm), \opcode{37} (Reg). Opcode \opcode{27} is undocumented.

	\textbf{Instruction Fields:}
	\begin{itemize}
		\item \textbf{Register:} destination \texttt{rD} (immediate form) or \texttt{rD, rS} (register form); any of the
		      16 registers, \texttt{0x0}--\texttt{0xF}.
		\item \textbf{Immediate:} 16-bit value (immediate form, bits 21--6). The register form has no immediate field.
		\item \textbf{AF (bits 5--4):} the public assembly syntax always encodes \texttt{00} (None), which preserves the
		      status flags. An encoding of this opcode with a different AF value, not reachable from assembly syntax,
		      leaves all four status flags architecturally undefined (\texttt{U U U U}).
		\item \textbf{Condition field (bits 3--0):} any of the 16 condition-code encodings in
		      Table~\ref{tab:condition-code-encoding}.
	\end{itemize}

	\textbf{Syntax:}
	\begin{lstlisting}
LOAD rD, #imm16                ; rD = imm16
LOAD rD, rS                    ; rD = rS
\end{lstlisting}

	\textbf{Operands:} Immediate form writes a 16-bit immediate value to \texttt{rD}. Register form copies \texttt{rS} to
	\texttt{rD}. No short-immediate, memory, or shift forms are public assembly syntax; the register-to-register form
	does not accept a shift definition (use \mnemonic{SHFT}). The indirect memory forms of \mnemonic{LOAD} are
	documented in Chapter~\ref{ch:memory-instructions}.

	\textbf{Operation:}
	\begin{verbatim}
destination = operand
; Flags NOT updated
\end{verbatim}

	\textbf{Description:}

	The ALU passes the second operand straight through to the destination; the first operand is ignored. This is the
	primary way to load constants or move data between registers. Opcode \texttt{0x27} has no public assembly syntax and is
	undocumented.

	\textbf{Flags:} Preserved.

	\textbf{Write-back:} If the condition is true, writes the loaded or copied value to \texttt{rD}.

	\textbf{Exceptions:} Does not access memory in these forms and raises no instruction-specific faults.

	\textbf{Condition field:} If the condition is false, no result is written and status flags are preserved.

	\textbf{Timing:} 6 cycles, plus instruction-fetch wait states.

	\textbf{Privilege:} Available in protected mode and supervisor mode, except when the register field names a
	privileged register (\reg{sr}, \reg{ah}, \reg{lh}, \reg{ph}, \reg{sh}); see Chapter~\ref{ch:exceptions}.

	\textbf{Example 13-10:}
	\begin{lstlisting}
; Load constant
LOAD r1, #1234             ; r1 = 1234

; Copy register
LOAD r2, r1                ; r2 = r1

; Load with arithmetic
LOAD r3, #1
ADDI r3, #0, LSL #8        ; r3 = (r3 << 8) + 0 = 256

; Zero a register
LOAD r4, #0                ; r4 = 0

; Load large constant (combine)
LOAD r5, #0x12
ADDI r5, #0, LSL #8        ; r5 = (r5 << 8) + 0 = 0x1200
ORRI r5, #0x34             ; r5 = 0x1234
\end{lstlisting}

	\textbf{Notes:}
	\begin{itemize}
		\item Does not update flags, unlike \mnemonic{ADD}
		\item Most efficient way to move data between registers
		\item For 16-bit constants, use immediate form
		\item To build shifted or composite constants, combine \mnemonic{LOAD} with ALU instructions such as \mnemonic{ADDI} or
		      \mnemonic{ORRI}
	\end{itemize}
\end{instructionbox}

% ============================================================================
\section{CMP -- Compare}

\begin{instructionbox}{CMPI / CMPR -- Compare Immediate / Compare Register}
	\textbf{Applicability:} SIRC-1, all revisions

	\textbf{Opcodes:} \opcode{0A} (Imm), \opcode{2A} (Short Imm), \opcode{3A} (Reg)

	\textbf{Instruction Fields:}
	\begin{itemize}
		\item \textbf{Register:} source \texttt{rS} (immediate and short-immediate forms) or \texttt{rS1, rS2} (register
		      form); any of the 16 registers, \texttt{0x0}--\texttt{0xF}. No destination register is written.
		\item \textbf{Immediate:} 16-bit value (immediate form, bits 21--6) or 8-bit value zero-extended to 16 bits
		      (short-immediate form, bits 21--14).
		\item \textbf{Shift type and count (short-immediate and register forms only):} \texttt{NUL}, \texttt{LSL},
		      \texttt{LSR}, \texttt{ASL}, \texttt{ASR}, \texttt{RTL}, \texttt{RTR}; an immediate count is
		      \texttt{0}--\texttt{15}; a register count is used in full (16 or more clamps to a full shift for \texttt{LSL}/\texttt{LSR}/\texttt{ASL}/\texttt{ASR}, and is taken modulo 16 for \texttt{RTL}/\texttt{RTR}) -- see Chapter~\ref{ch:shift-operations}.
		\item \textbf{AF (bits 5--4):} \texttt{00} = None, \texttt{01} = Alu (default), \texttt{10} = Shift, \texttt{11} =
		      Reserved (status register unchanged).
		\item \textbf{Condition field (bits 3--0):} any of the 16 condition-code encodings in
		      Table~\ref{tab:condition-code-encoding}.
	\end{itemize}

	\textbf{Syntax:}
	\begin{lstlisting}
CMPI rS, #imm16                ; Compare rS with imm16
CMPI rS, #imm8, shift          ; Compare shift(rS) with imm8
CMPR rS1, rS2                  ; Compare rS1 with rS2
CMPR rS1, rS2, shift           ; Compare shift(rS1) with rS2
\end{lstlisting}

	\textbf{Operands:} Same source operand forms as \mnemonic{SUB}. The first operand is compared with the second operand;
	no destination register is written.

	\textbf{Operation:}
	\begin{verbatim}
result = operand1 - operand2
; Result is discarded, only flags updated
SR.N = result[15]
SR.Z = (result == 0)
SR.C = borrow
SR.V = signed_overflow
\end{verbatim}

	\textbf{Description:}

	Performs subtraction but discards the result, only updating the status flags. In short immediate format with shift, the
	source register is shifted before comparing with the immediate. This is the primary instruction for implementing
	conditional execution and branches. The flags can be tested with condition codes to determine the relationship between
	the operands.

	\textbf{Flags:} N Z C V

	\textbf{Write-back:} None. The subtraction result is discarded after status flags are updated.

	\textbf{Exceptions:} Does not access memory and raises no instruction-specific faults.

	\textbf{Condition field:} If the condition is false, no status flags are updated.

	\textbf{Timing:} 6 cycles, plus instruction-fetch wait states.

	\textbf{Privilege:} Available in protected mode and supervisor mode, except when the register field names a
	privileged register (\reg{sr}, \reg{ah}, \reg{lh}, \reg{ph}, \reg{sh}); see Chapter~\ref{ch:exceptions}. The
	privilege check is made on the register field even though this instruction writes no register, so
	\texttt{CMPI sr, \#0} in protected mode raises a privilege-violation fault.

	\textbf{Example 13-11:}
	\begin{lstlisting}
; Test if r1 is zero
CMPI r1, #0
BRAN|== @is_zero           ; Branch if r1 == 0

; Test if r1 > r2 (unsigned)
CMPR r1, r2
BRAN|HI @r1_greater

; Test if r1 >= r2 (signed)
CMPR r1, r2
BRAN|>= @r1_greater_or_equal

; Range check: if (r1 >= 10 && r1 < 20)
CMPI r1, #10
BRAN|<< @out_of_range      ; Branch if r1 < 10
CMPI r1, #20
BRAN|>= @out_of_range      ; Branch if r1 >= 20
; r1 is in range [10, 20)
\end{lstlisting}

	\textbf{Notes:}
	\begin{itemize}
		\item Does not modify any register, only flags
		\item Use with conditional instructions or branches
		\item Remember: \texttt{CMP A, B} computes A - B
		\item For signed comparisons, use \texttt{>=}, \texttt{<<}, \texttt{>>}, \texttt{<=}
		\item For unsigned comparisons, use \texttt{HI}, \texttt{LO}, \texttt{CS}, \texttt{CC}
	\end{itemize}
\end{instructionbox}

% ============================================================================
\section{TSA -- Test AND}

\begin{instructionbox}{TSAI / TSAR -- Test AND Immediate / Test AND Register}
	\textbf{Applicability:} SIRC-1, all revisions

	\textbf{Opcodes:} \opcode{0C} (Imm), \opcode{2C} (Short Imm), \opcode{3C} (Reg)

	\textbf{Instruction Fields:}
	\begin{itemize}
		\item \textbf{Register:} source \texttt{rS} (immediate and short-immediate forms) or \texttt{rS1, rS2} (register
		      form); any of the 16 registers, \texttt{0x0}--\texttt{0xF}. No destination register is written.
		\item \textbf{Immediate:} 16-bit value (immediate form, bits 21--6) or 8-bit value zero-extended to 16 bits
		      (short-immediate form, bits 21--14).
		\item \textbf{Shift type and count (short-immediate and register forms only):} \texttt{NUL}, \texttt{LSL},
		      \texttt{LSR}, \texttt{ASL}, \texttt{ASR}, \texttt{RTL}, \texttt{RTR}; an immediate count is
		      \texttt{0}--\texttt{15}; a register count is used in full (16 or more clamps to a full shift for \texttt{LSL}/\texttt{LSR}/\texttt{ASL}/\texttt{ASR}, and is taken modulo 16 for \texttt{RTL}/\texttt{RTR}) -- see Chapter~\ref{ch:shift-operations}.
		\item \textbf{AF (bits 5--4):} \texttt{00} = None, \texttt{01} = Alu (default), \texttt{10} = Shift, \texttt{11} =
		      Reserved (status register unchanged).
		\item \textbf{Condition field (bits 3--0):} any of the 16 condition-code encodings in
		      Table~\ref{tab:condition-code-encoding}.
	\end{itemize}

	\textbf{Syntax:}
	\begin{lstlisting}
TSAI rS, #imm16                ; Test rS & imm16
TSAI rS, #imm8, shift          ; Test shift(rS) & imm8
TSAR rS1, rS2                  ; Test rS1 & rS2
TSAR rS1, rS2, shift           ; Test shift(rS1) & rS2
\end{lstlisting}

	\textbf{Operands:} Same source operand forms as \mnemonic{AND}. The result is used only for status flag updates; no
	destination register is written.

	\textbf{Operation:}
	\begin{verbatim}
result = operand1 AND operand2
; Result is discarded, only flags updated
SR.N = result[15]
SR.Z = (result == 0)
SR.C = 0
SR.V = 0
\end{verbatim}

	\textbf{Description:}

	Performs bitwise AND but discards the result, only updating flags. In short immediate format with shift, the source
	register is shifted before ANDing with the immediate. Used to test if specific bits are set without modifying the
	register. The zero flag indicates whether any tested bits were set.

	\textbf{Flags:} N Z (C and V cleared)

	\textbf{Write-back:} None. The AND result is discarded after status flags are updated.

	\textbf{Exceptions:} Does not access memory and raises no instruction-specific faults.

	\textbf{Condition field:} If the condition is false, no status flags are updated.

	\textbf{Timing:} 6 cycles, plus instruction-fetch wait states.

	\textbf{Privilege:} Available in protected mode and supervisor mode, except when the register field names a
	privileged register (\reg{sr}, \reg{ah}, \reg{lh}, \reg{ph}, \reg{sh}); see Chapter~\ref{ch:exceptions}. The
	privilege check is made on the register field even though this instruction writes no register, so
	\texttt{TSAI sr, \#0} in protected mode raises a privilege-violation fault.

	\textbf{Example 13-12:}
	\begin{lstlisting}
; Test if bit 5 is set
TSAI r1, #0x0020           ; Test r1 & 0x20
BRAN|!= @bit5_set          ; Branch if bit 5 was set

; Test if any high byte bits are set
TSAI r2, #0xFF00
BRAN|!= @has_upper_bits

; Test if r1 and r2 have any bits in common
TSAR r1, r2
BRAN|== @no_common_bits    ; Branch if no bits in common

; Test multiple bits
TSAI r3, #0xF000           ; Test high nibble
BRAN|== @upper_clear       ; Branch if all tested bits clear
\end{lstlisting}

	\textbf{Notes:}
	\begin{itemize}
		\item Does not modify any register, only flags
		\item Z=1 means all tested bits were 0
		\item Z=0 means at least one tested bit was 1
		\item Use for bit testing without side effects
		\item More efficient than AND + CMP for simple tests
	\end{itemize}
\end{instructionbox}

% ============================================================================
\section{TSX -- Test XOR}

\begin{instructionbox}{TSXI / TSXR -- Test XOR Immediate / Test XOR Register}
	\textbf{Applicability:} SIRC-1, all revisions

	\textbf{Opcodes:} \opcode{0E} (Imm), \opcode{2E} (Short Imm), \opcode{3E} (Reg)

	\textbf{Instruction Fields:}
	\begin{itemize}
		\item \textbf{Register:} source \texttt{rS} (immediate and short-immediate forms) or \texttt{rS1, rS2} (register
		      form); any of the 16 registers, \texttt{0x0}--\texttt{0xF}. No destination register is written.
		\item \textbf{Immediate:} 16-bit value (immediate form, bits 21--6) or 8-bit value zero-extended to 16 bits
		      (short-immediate form, bits 21--14).
		\item \textbf{Shift type and count (short-immediate and register forms only):} \texttt{NUL}, \texttt{LSL},
		      \texttt{LSR}, \texttt{ASL}, \texttt{ASR}, \texttt{RTL}, \texttt{RTR}; an immediate count is
		      \texttt{0}--\texttt{15}; a register count is used in full (16 or more clamps to a full shift for \texttt{LSL}/\texttt{LSR}/\texttt{ASL}/\texttt{ASR}, and is taken modulo 16 for \texttt{RTL}/\texttt{RTR}) -- see Chapter~\ref{ch:shift-operations}.
		\item \textbf{AF (bits 5--4):} \texttt{00} = None, \texttt{01} = Alu (default), \texttt{10} = Shift, \texttt{11} =
		      Reserved (status register unchanged).
		\item \textbf{Condition field (bits 3--0):} any of the 16 condition-code encodings in
		      Table~\ref{tab:condition-code-encoding}.
	\end{itemize}

	\textbf{Syntax:}
	\begin{lstlisting}
TSXI rS, #imm16                ; Test rS ^ imm16
TSXI rS, #imm8, shift          ; Test shift(rS) ^ imm8
TSXR rS1, rS2                  ; Test rS1 ^ rS2
TSXR rS1, rS2, shift           ; Test shift(rS1) ^ rS2
\end{lstlisting}

	\textbf{Operands:} Same source operand forms as \mnemonic{XOR}. The result is used only for status flag updates; no
	destination register is written.

	\textbf{Operation:}
	\begin{verbatim}
result = operand1 XOR operand2
; Result is discarded, only flags updated
SR.N = result[15]
SR.Z = (result == 0)
SR.C = 0
SR.V = 0
\end{verbatim}

	\textbf{Description:}

	Performs bitwise XOR but discards the result, only updating flags. In short immediate format with shift, the source
	register is shifted before XORing with the immediate. Used to test if bits match a specific pattern or to compare
	values for equality. Z=1 indicates the values are identical.

	\textbf{Flags:} N Z (C and V cleared)

	\textbf{Write-back:} None. The XOR result is discarded after status flags are updated.

	\textbf{Exceptions:} Does not access memory and raises no instruction-specific faults.

	\textbf{Condition field:} If the condition is false, no status flags are updated.

	\textbf{Timing:} 6 cycles, plus instruction-fetch wait states.

	\textbf{Privilege:} Available in protected mode and supervisor mode, except when the register field names a
	privileged register (\reg{sr}, \reg{ah}, \reg{lh}, \reg{ph}, \reg{sh}); see Chapter~\ref{ch:exceptions}. The
	privilege check is made on the register field even though this instruction writes no register, so
	\texttt{TSXI sr, \#0} in protected mode raises a privilege-violation fault.

	\textbf{Example 13-13:}
	\begin{lstlisting}
; Test if r1 equals specific value
TSXI r1, #0x1234           ; Test r1 ^ 0x1234
BRAN|== @is_0x1234         ; Branch if r1 == 0x1234

; Test if two registers are equal
TSXR r1, r2                ; Test r1 ^ r2
BRAN|== @registers_equal   ; Branch if r1 == r2

; Test if bit pattern matches
ANDI r3, #0x00FF           ; Mask to low byte
TSXI r3, #0x00AA           ; Test if low byte == 0xAA
BRAN|== @pattern_match
\end{lstlisting}

	\textbf{Notes:}
	\begin{itemize}
		\item Does not modify any register, only flags
		\item Z=1 means values are identical (all bits match)
		\item Can be used as equality test (alternative to CMP for == only)
		\item More efficient than XOR + CMP for simple equality tests
	\end{itemize}
\end{instructionbox}

\section{Summary}

The ALU instructions form the computational core of the SIRCIS instruction set. Key points:

\begin{itemize}
	\item All ALU operations work only on registers (load/store architecture)
	\item Three format variants provide flexibility (immediate, short+shift, register)
	\item Test variants (CMP, TSA, TSX) allow non-destructive testing
	\item Shift operations can be combined with ALU ops for powerful single-instruction operations
	\item Consistent flag behavior makes conditional execution predictable
\end{itemize}
